\chapter{Koleksi Prompt Engineering Terbaik}
\label{chap:prompts}

Appendix ini berisi kumpulan prompt siap pakai yang telah dioptimalkan untuk berbagai kebutuhan. Setiap prompt dirancang menggunakan prinsip "Context, Instruction, Output Format".

\section{Coding \& Software Development}

\subsection{1. Menjelaskan Kode Kompleks}
\begin{promptbox}
Saya memiliki potongan kode Python berikut yang sulit saya pahami. Bertindaklah sebagai Senior Python Developer.
1. Jelaskan logika kode ini baris demi baris.
2. Identifikasi potensi bug atau inefisiensi (Big O notation).
3. Tulis ulang kode agar lebih Pythonic dan efisien.

[PASTE KODE DI SINI]
\end{promptbox}

\subsection{2. Membuat Unit Test (Pytest)}
\begin{promptbox}
Tulis unit test case lengkap menggunakan framework `pytest` untuk fungsi berikut. Pastikan untuk mencakup:
- Happy path (input valid).
- Edge cases (input kosong, nilai negatif, tipe data salah).
- Mocking jika ada panggilan API eksternal.

[PASTE FUNGSI DI SINI]
\end{promptbox}

\subsection{3. Mengonversi Bahasa Pemrograman}
\begin{promptbox}
Konversi kode berikut dari Java ke Python 3.10+.
- Gunakan type hinting Python.
- Ganti loop for-loop gaya Java dengan list comprehension jika memungkinkan.
- Pertahankan nama variabel agar tetap konsisten.

[PASTE KODE JAVA DI SINI]
\end{promptbox}

\subsection{4. Debugging Error}
\begin{promptbox}
Saya mendapatkan error berikut saat menjalankan aplikasi Node.js saya.
Error: [PASTE ERROR LOG]

Konteks kode:
[PASTE SNIPPET KODE]

Analisis penyebab error ini dan berikan 3 solusi potensial, diurutkan dari yang paling mungkin berhasil.
\end{promptbox}

\subsection{5. Menulis Dokumentasi API}
\begin{promptbox}
Buat dokumentasi API dalam format OpenAPI (Swagger) YAML untuk endpoint berikut.
Endpoint: POST /api/v1/users/register
Input JSON: {username, password, email}
Output JSON: {id, created_at, token}
Error Responses: 400 (Bad Request), 409 (Username Taken).

Sertakan deskripsi detail untuk setiap field.
\end{promptbox}

\section{Content Writing & Copywriting}

\subsection{6. Membuat Outline Artikel Blog SEO}
\begin{promptbox}
Bertindaklah sebagai SEO Specialist. Buat outline artikel blog yang komprehensif tentang topik "[TOPIK]".
Target audiens: [AUDIENS].
Keyword utama: [KEYWORD].
Struktur yang diinginkan:
- Judul H1 yang menarik (clickbait tapi jujur).
- Introduction dengan hook.
- 5-7 Subheading (H2) yang menjawab search intent.
- FAQ section (H2).
- Conclusion dengan Call to Action (CTA).
\end{promptbox}

\subsection{7. Menulis Thread Twitter/X Viral}
\begin{promptbox}
Ubah artikel/teks berikut menjadi Thread Twitter yang viral.
- Tweet pertama harus berupa "hook" yang memancing rasa ingin tahu.
- Gunakan bahasa yang santai, ringkas, dan punchy.
- Setiap tweet maksimal 280 karakter.
- Tambahkan emoji yang relevan.
- Tweet terakhir adalah ringkasan/CTA.

[PASTE TEKS SUMBER]
\end{promptbox}

\subsection{8. Memperbaiki Tata Bahasa & Gaya (Editing)}
\begin{promptbox}
Bertindaklah sebagai editor profesional. Perbaiki tulisan berikut agar lebih:
- Ringkas (hapus kata-kata filler).
- Persuasif.
- Menggunakan nada suara yang [FORMAL/SANTAI/HUMORIS].
- Perbaiki kesalahan tata bahasa (EYD).

Teks asli:
[PASTE TEKS]
\end{promptbox}

\subsection{9. Membuat Skrip Video YouTube/TikTok}
\begin{promptbox}
Buat naskah video TikTok durasi 60 detik tentang "[TOPIK]".
Format tabel dengan 2 kolom:
1. Visual (apa yang terlihat di layar, B-roll, teks overlay).
2. Audio (apa yang diucapkan narator).

Gunakan hook visual di 3 detik pertama untuk menahan retensi penonton.
\end{promptbox}

\subsection{10. Cold Email untuk Penjualan}
\begin{promptbox}
Tulis email penawaran (cold email) ke CEO perusahaan [JENIS PERUSAHAAN] untuk menawarkan jasa [JASA ANDA].
- Subject line: Pendek, personal, tidak terlihat seperti spam.
- Paragraf 1: Relate dengan masalah mereka (pain point).
- Paragraf 2: Tawarkan solusi/value proposition.
- Paragraf 3: Social proof (kami telah membantu X perusahaan).
- CTA: Ajakan untuk call singkat 15 menit.
\end{promptbox}

\section{Bisnis & Produktivitas}

\subsection{11. Meringkas Dokumen Panjang}
\begin{promptbox}
Ringkas laporan keuangan/dokumen ini ke dalam format "Executive Summary" satu halaman.
Sertakan:
- Poin-poin utama (bullet points).
- Data/angka kunci.
- Kesimpulan strategis.
- Rekomendasi tindakan.

[PASTE DOKUMEN]
\end{promptbox}

\subsection{12. Analisis SWOT}
\begin{promptbox}
Lakukan analisis SWOT (Strengths, Weaknesses, Opportunities, Threats) untuk ide bisnis berikut:
"[IDE BISNIS ANDA]"
Asumsikan pasar target adalah [LOKASI/DEMOGRAFI]. Berikan strategi mitigasi untuk setiap Weakness dan Threat.
\end{promptbox}

\subsection{13. Membuat Rencana Proyek (Gantt Chart text)}
\begin{promptbox}
Saya ingin meluncurkan [PROYEK, misal: Website E-commerce] dalam waktu 3 bulan.
Buat rencana kerja mingguan (week-by-week breakdown).
Format:
- Minggu 1-2: Fase Perencanaan
- Minggu 3-6: Fase Pengembangan
- dst.
Sertakan deliverable konkret di setiap fase.
\end{promptbox}

\subsection{14. Persiapan Negosiasi Gaji}
\begin{promptbox}
Saya akan melakukan negosiasi gaji untuk posisi [POSISI] di [KOTA]. Tawaran saat ini adalah [ANGKA], tapi target saya adalah [TARGET].
- Berikan skrip percakapan untuk menanggapi tawaran awal.
- Daftar argumen yang bisa saya gunakan (pengalaman, inflasi, nilai pasar).
- Cara menolak dengan sopan jika angka tidak masuk.
\end{promptbox}

\subsection{15. Ide Brainstorming Produk}
\begin{promptbox}
Gunakan teknik SCAMPER (Substitute, Combine, Adapt, Modify, Put to another use, Eliminate, Reverse) untuk menghasilkan 7 variasi inovatif dari produk: [NAMA PRODUK, misal: Payung].
\end{promptbox}

\section{Pendidikan & Belajar}

\subsection{16. Tutor Pribadi (Socratic Method)}
\begin{promptbox}
Saya ingin belajar tentang [TOPIK, misal: Black Holes].
Jangan langsung berikan kuliah panjang. Bertindaklah sebagai tutor Socrates.
Ajukan satu pertanyaan sederhana kepada saya untuk menguji pemahaman awal saya.
Tunggu jawaban saya, lalu berikan feedback dan ajukan pertanyaan berikutnya yang lebih dalam. Terus bimbing saya sampai saya memahami konsep intinya.
\end{promptbox}

\subsection{17. Membuat Soal Kuis}
\begin{promptbox}
Buat 5 soal pilihan ganda tentang materi [MATERI].
- Tingkat kesulitan: [MUDAH/SEDANG/SULIT].
- Setiap soal harus memiliki 4 opsi (A, B, C, D).
- Sertakan kunci jawaban dan penjelasan singkat di akhir (setelah semua soal).
\end{promptbox}

\subsection{18. Menjelaskan ke Anak 5 Tahun (ELI5)}
\begin{promptbox}
Jelaskan konsep [KONSEP SULIT, misal: Inflasi] seolah-olah saya adalah anak usia 5 tahun. Gunakan analogi sederhana (seperti mainan atau permen) dan hindari jargon teknis.
\end{promptbox}

\subsection{19. Membuat Rencana Belajar (Sylabus)}
\begin{promptbox}
Buat kurikulum belajar mandiri selama 4 minggu untuk menguasai skill [SKILL, misal: Video Editing di Premiere Pro].
- Alokasi waktu: 1 jam per hari.
- Hari 1-5: Materi & Latihan.
- Hari 6: Proyek Mingguan.
- Hari 7: Review.
Sertakan topik spesifik yang harus dipelajari setiap hari.
\end{promptbox}

\subsection{20. Flashcard Generator}
\begin{promptbox}
Buat tabel 2 kolom untuk Anki/Flashcard dari teks berikut.
Kolom 1: Istilah/Konsep (Depan Kartu).
Kolom 2: Definisi/Penjelasan Singkat (Belakang Kartu).

[PASTE TEKS]
\end{promptbox}

\section{Data Analysis & Excel}

\subsection{21. Membuat Rumus Excel/Google Sheets}
\begin{promptbox}
Saya punya data di Kolom A (Tanggal Lahir) dan ingin menghitung Umur di Kolom B.
Tuliskan rumus Excel yang akurat (memperhitungkan tahun kabisat dan tanggal hari ini).
Juga, bagaimana cara memformat sel agar tampil "XX Tahun, YY Bulan"?
\end{promptbox}

\subsection{22. Analisis Sentimen Teks}
\begin{promptbox}
Analisis sentimen dari ulasan pelanggan berikut. Kategorikan menjadi: Positif, Netral, atau Negatif.
Buat dalam format tabel CSV:
Ulasan | Sentimen | Kata Kunci Emosi

Ulasan 1: ...
Ulasan 2: ...
\end{promptbox}

\subsection{23. Interpretasi Data SQL}
\begin{promptbox}
Tulis query SQL untuk:
"Menampilkan 5 pelanggan teratas yang memiliki total belanja tertinggi di bulan Januari 2024, beserta total belanjanya."
Nama tabel: `orders` dan `customers`.
Kolom: `customer_id`, `order_date`, `amount`.
\end{promptbox}

\subsection{24. Regex Generator}
\begin{promptbox}
Buat Regular Expression (Regex) untuk mengekstrak semua alamat email dari teks yang berantakan.
Regex harus bisa menangkap format:
- name@domain.com
- name.surname@sub.domain.co.id
Jelaskan cara kerja pola regex tersebut.
\end{promptbox}

\subsection{25. Data Cleaning Python}
\begin{promptbox}
Saya punya DataFrame pandas `df`. Kolom 'price' bertipe object karena ada simbol mata uang ('Rp 15.000').
Tulis kode Python untuk:
1. Menghapus 'Rp ' dan titik.
2. Mengubah tipe data menjadi integer.
3. Menangani nilai NaN dengan mengisinya dengan rata-rata.
\end{promptbox}

\section{Marketing & Sales}

\subsection{26. Membuat Persona Pelanggan (Buyer Persona)}
\begin{promptbox}
Buat profil Buyer Persona detail untuk produk [PRODUK, misal: Kopi Susu Kekinian].
Sertakan:
- Nama Fiktif & Demografi (Umur, Pekerjaan).
- Goals & Values.
- Pain Points & Frustrations.
- Media yang dikonsumsi (Instagram/TikTok/LinkedIn).
- Objections (kenapa mereka ragu membeli).
\end{promptbox}

\subsection{27. Ide Kampanye Marketing}
\begin{promptbox}
Berikan 5 ide kampanye marketing unik dan berbiaya rendah (guerrilla marketing) untuk mempromosikan [BISNIS ANDA].
Fokus pada viralitas di media sosial.
\end{promptbox}

\subsection{28. Menulis Ad Copy (Facebook/Instagram Ads)}
\begin{promptbox}
Tulis 3 variasi teks iklan Facebook untuk produk [PRODUK].
Gunakan kerangka AIDA (Attention, Interest, Desire, Action).
- Variasi 1: Pendek & Punchy.
- Variasi 2: Storytelling (Testimonial).
- Variasi 3: FOMO (Fear of Missing Out).
\end{promptbox}

\subsection{29. Email Newsletter}
\begin{promptbox}
Buat naskah newsletter mingguan untuk audiens [AUDIENS].
Tema minggu ini: [TEMA].
Struktur:
- Salam hangat.
- Cerita pendek/anekdot terkait tema.
- 3 Tips praktis.
- Link ke konten terbaru.
- Penutup.
\end{promptbox}

\subsection{30. Review Respons}
\begin{promptbox}
Tulis balasan yang sopan dan profesional untuk review negatif pelanggan berikut di Google Maps.
Review: "[Review Negatif Pelanggan]"
Tujuan: Minta maaf atas ketidaknyamanan, tawarkan solusi, dan ajak diskusi lebih lanjut secara privat (DM/Email) agar tidak berdebat di publik.
\end{promptbox}

\section{Prompting untuk Gambar (Midjourney/DALL-E)}

\subsection{31. Fotografi Realistis}
\begin{promptbox}
Hyper-realistic portrait of an elderly Indonesian fisherman, weathered face, deep wrinkles, looking at the horizon, sunset lighting, golden hour, cinematic bokeh, 85mm lens, f/1.8, highly detailed texture, 8k resolution --ar 16:9 --v 6.0
\end{promptbox}

\subsection{32. Logo Desain}
\begin{promptbox}
Minimalist vector logo for a coffee shop named "Kopi Senja", icon combining a coffee bean and a setting sun, flat design, simple lines, orange and dark brown color palette, white background --no text --v 6.0
\end{promptbox}

\subsection{33. Ilustrasi 3D (Pixar Style)}
\begin{promptbox}
3D cute isometric bedroom interior, cozy lighting, messy desk with computer, cat sleeping on chair, plants, Pixar animation style, soft render, octane render, blender 3d, pastel colors --ar 1:1
\end{promptbox}

\subsection{34. Cyberpunk City}
\begin{promptbox}
Futuristic Jakarta city skyline at night, cyberpunk style, neon lights, rain reflections on wet asphalt, flying cars, holograms, towering skyscrapers, dystopian atmosphere, cinematic composition --ar 21:9
\end{promptbox}

\subsection{35. UI/UX Design Inspiration}
\begin{promptbox}
Mobile app UI design for a travel booking application, clean and modern interface, trending on Dribbble, Figma style, blue and white color scheme, high fidelity, user friendly layout --ar 9:16
\end{promptbox}

\section{Personal Assistant & Life Hacking}

\subsection{36. Perencanaan Perjalanan (Itinerary)}
\begin{promptbox}
Buatkan itinerary perjalanan 5 hari 4 malam ke Jepang (Tokyo & Kyoto) untuk keluarga dengan 2 anak kecil.
- Budget: Menengah.
- Minat: Budaya, Kuliner, dan Taman Bermain.
- Sertakan opsi transportasi antar kota.
- Jangan terlalu padat, berikan waktu istirahat.
\end{promptbox}

\subsection{37. Meal Plan Mingguan (Diet)}
\begin{promptbox}
Buatkan rencana makan (meal plan) mingguan yang sehat dan murah (budget Rp 500rb/minggu).
- Tujuan: Menurunkan berat badan (defisit kalori).
- Preferensi: Makanan Indonesia, mudah dimasak (maks 30 menit).
- Alergi: Tidak ada kacang.
Sertakan daftar belanjaan (shopping list) lengkap.
\end{promptbox}

\subsection{38. Rekomendasi Buku/Film}
\begin{promptbox}
Saya sangat menyukai film "Interstellar" dan "Inception".
Berikan 5 rekomendasi film lain yang memiliki tema "mind-bending sci-fi" atau "time dilation" serupa.
Jelaskan secara singkat mengapa saya akan menyukainya (tanpa spoiler).
\end{promptbox}

\subsection{39. Format Email Profesional (Resign)}
\begin{promptbox}
Tulis surat pengunduran diri (resignation letter) yang sopan dan profesional.
- Alasan: Mendapat tawaran yang lebih baik.
- Notice period: 1 bulan.
- Nada: Berterima kasih atas kesempatan yang diberikan dan ingin menjaga hubungan baik.
\end{promptbox}

\subsection{40. Simulasi Debat (Devil's Advocate)}
\begin{promptbox}
Saya berpendapat bahwa "Remote Working (WFH) lebih baik daripada WFO".
Bertindaklah sebagai lawan debat yang cerdas. Berikan 3 argumen kontra yang kuat dan berbasis data mengapa WFO (Work From Office) sebenarnya lebih unggul untuk inovasi dan budaya perusahaan.
\end{promptbox}

\section{Creative Writing \& Storytelling}

\subsection{41. Membangun Dunia (World Building)}
\begin{promptbox}
Saya sedang menulis novel fantasi. Bantu saya membangun sistem sihir (magic system) yang unik.
- Aturan dasar: Sihir membutuhkan "pengorbanan" memori.
- Batasan: Tidak bisa menghidupkan orang mati.
- Konsekuensi sosial: Pengguna sihir dikucilkan atau dipuja?
Berikan 3 konsep sistem sihir berbeda berdasarkan premis ini.
\end{promptbox}

\subsection{42. Pengembangan Karakter}
\begin{promptbox}
Buat lembar karakter (character sheet) untuk antagonis utama cerita thriller psikologis.
- Nama: [NAMA]
- Motivasi: Bukan uang atau kekuasaan, tapi "menyelamatkan" orang dengan cara yang salah.
- Latar belakang traumatis.
- Kebiasaan kecil (quirk) yang mengganggu.
- Hubungan dengan protagonis.
\end{promptbox}

\subsection{43. Plot Twist Generator}
\begin{promptbox}
Cerita saya adalah tentang detektif yang menyelidiki kasus pembunuhan berantai di stasiun luar angkasa.
Berikan 5 ide plot twist yang mencengangkan untuk ending cerita.
Salah satu twist harus mengubah genre cerita dari Misteri menjadi Horor Kosmik.
\end{promptbox}

\subsection{44. Menulis Puisi}
\begin{promptbox}
Tulis puisi bebas (free verse) tentang "Hujan di Jakarta" yang melankolis tapi penuh harapan.
Gunakan citraan (imagery) tentang kemacetan, lampu neon, dan aroma tanah basah (petrichor).
Hindari rima yang kaku.
\end{promptbox}

\subsection{45. Skenario Dialog}
\begin{promptbox}
Tulis dialog intens antara dua sahabat yang sudah lama tidak bertemu.
Konteks: Sahabat A baru saja keluar dari penjara karena kesalahan Sahabat B, tapi Sahabat B tidak tahu bahwa A tahu.
Subtext: Ketegangan, rasa bersalah, dan kemarahan yang ditahan.
Setting: Sebuah warung kopi yang bising.
\end{promptbox}

\begin{tipbox}
    \textbf{Pro Tip:} Simpan prompt terbaik Anda dalam dokumen terpisah (seperti Notion atau Obsidian) agar bisa digunakan kembali. Prompt adalah aset digital Anda.
\end{tipbox}
